\section{Kernel exclusion, paired heat transport and arithmetic approximation (v1.54)}
\label{sec:v154-new-shapes}
This checkpoint tests three different intermediate mechanisms. It proves an
exact complete-domain nullvector reduction and a scoped regularity countermodel,
finite heat-zero transport identities with an explicitly conditional comparison
obstruction, and an unconditional arithmetic block-gain inequality. The missing
inputs are, respectively, arithmetic injectivity, suitable full heat transport,
and asymptotic residual correlation. None is supplied by changing coordinates.
The countermodels and implications below do not produce a negative Weil
direction. No new window or tail metric is computed; G2 and RH remain open.

\subsection*{The complete nullvector equation without an $H^1$ assumption (NS-51)}

\paragraph{Classification and scope.}
The next statements are exact identities and proved implications, with a
named open arithmetic injectivity input. They concern the complete
fixed-window operator, not a finite head. The new regularity countermodel
rules out one generic inference, not injectivity of the arithmetic
operator. No new window or tail metric is computed. G2 and RH remain open.

Put $I_a=(-a,a)$, and let $E_a$ denote zero extension from $I_a$.
Write $V_a=\mathcal D(q_a)$ for the complete physical form domain:
\[
 V_a=\left\{f\in L^2(I_a):
  \int_{\mathbb R}\log(2+|\xi|)|\widehat{E_af}(\xi)|^2\,d\xi<\infty\right\}.
\]
Let $A_a$ be its self-adjoint form operator. These are the unshifted,
complete objects of Proposition~\ref{prop:v118-log-dirichlet}; no source
projection is imposed. Let $\mathfrak g$ be Suzuki's real even screw
function, normalized by $\mathfrak g(0)=0$ and, for $t>0$,
\begin{align}
 \mathfrak g(t)={}&-4(e^{t/2}+e^{-t/2}-2)
  +\sum_{n<e^t}\frac{\Lambda(n)}{\sqrt n}(t-\log n)
  -\frac t2\bigl(\psi(1/4)-\log\pi\bigr)\notag\\
 &-\frac14\left\{\Phi(1,2,1/4)
             -e^{-t/2}\Phi(e^{-2t},2,1/4)\right\}.
 \label{eq:ns51-screw}
\end{align}
Here $\Phi(z,s,v)=\sum_{j\ge0}z^j/(j+v)^s$ for $|z|<1$,
with its convergent value at $(1,2,1/4)$. Including a prime-power
endpoint instead makes no difference in this formula, since its ramp
vanishes at that endpoint. Suzuki's distributional identity is
$A_a\varphi=(-\mathfrak g''*E_a\varphi)|_{I_a}$ for
$\varphi\in C_c^\infty(I_a)$~\cite[Sections 2.5 and 3]{Suzuki2026}.
All convolutions below are local convolutions with a compactly supported
function or distribution; no global temperedness of $\mathfrak g$ is
assumed.

\begin{proposition}[Affine screw potential criterion on the complete domain]
\label{prop:ns51-affine-potential}
For $f\in L^2(I_a)$ set
\[
 F_f(x)=\int_{-a}^a\mathfrak g(x-y)f(y)\,dy,\qquad
 U_f(x)=\int_{-a}^a\mathfrak g'(x-y)f(y)\,dy.
\]
Then $F_f\in C^1(\mathbb R)$ and $F_f'=U_f\in C(\mathbb R)$.
For $f\in V_a$ the following conditions are equivalent:
\begin{enumerate}[label=(\roman*)]
 \item $f\in\mathcal D(A_a)$ and $A_af=0$;
 \item $F_f$ is affine on $I_a$;
 \item $U_f$ is constant on $I_a$.
\end{enumerate}
Let $\Pi_{\rm aff}$ be the ordinary $L^2(I_a)$ projection onto
$\operatorname{span}\{1,x\}$ and $P_0$ the projection onto mean-zero
functions. The operators
\[
 \mathcal C_a f=(I-\Pi_{\rm aff})F_f|_{I_a},\qquad
 \mathcal T_a f=P_0 U_f|_{I_a}
\]
are compact on $L^2(I_a)$, and
\begin{equation}\label{eq:ns51-kernel-reduction}
 \ker A_a=V_a\cap\ker\mathcal C_a=V_a\cap\ker\mathcal T_a.
\end{equation}
For even $f$ the affine potential is constant and $U_f=0$; for odd
$f$ it is $b x$ and $U_f=b$. The constant $b$ must not be discarded.
\end{proposition}
\begin{proof}
The screw function is locally absolutely continuous. Its derivative
has only a logarithmic singularity at zero and finitely many bounded
jumps on every compact interval. Indeed differentiation of
\eqref{eq:ns51-screw} gives, almost everywhere for $t>0$,
\begin{equation}\label{eq:ns51-screw-derivative}
 \mathfrak g'(t)=-4\sinh(t/2)
 +\sum_{n<e^t}\frac{\Lambda(n)}{\sqrt n}
 -\frac12(\psi(1/4)-\log\pi)
 -\frac12e^{-t/2}\Phi(e^{-2t},1,1/4),
\end{equation}
and odd reflection gives its value for $t<0$.
In particular $\mathfrak g'\in L^2_{\rm loc}(\mathbb R)$.
Translation continuity in local $L^2$ and Cauchy--Schwarz show that
$U_f$ is continuous, locally uniformly in $x$. The distributional
identity $F_f'=U_f$ then proves $F_f\in C^1$.

For $f\in V_a$ and $\varphi\in C_c^\infty(I_a)$, the core identity
and symmetry give, in the first-slot-linear convention,
\[
 q_a(f,\varphi)=\langle f,A_a\varphi\rangle
 =-\int_{I_a}F_f(x)\overline{\varphi''(x)}\,dx
 =\int_{I_a}U_f(x)\overline{\varphi'(x)}\,dx.
\]
A nullvector therefore has $F_f''=0$ as a distribution on $I_a$,
which is equivalent to the displayed affine and constant conditions.
Conversely these conditions annihilate $q_a(f,\varphi)$ on the smooth
form core. The form norm, after a fixed-window positive shift, makes
$q_a(f,\cdot)$ continuous. Core density therefore gives
$q_a(f,v)=0$ for every $v\in V_a$. The representation theorem for
closed forms implies $f\in\mathcal D(A_a)$ and $A_af=0$.
The kernels $\mathfrak g(x-y)$ and $\mathfrak g'(x-y)$ are square
integrable on $I_a^2$; for the latter the squared kernel norm is
\[
 \int_{-2a}^{2a}(2a-|t|)|\mathfrak g'(t)|^2\,dt<\infty.
\]
Thus their integral operators are Hilbert--Schmidt, and bounded
projections preserve compactness. Evenness of $\mathfrak g$ gives
the two parity conclusions.
\end{proof}

\paragraph{What changes in the derivative reduction.}
For $f\in H_0^1(I_a)$ one can write the familiar equation
$G_a Df=0$ with $D=i\,d/dx$ and the ordinary mean-zero screw operator
$G_a$. For general $f\in V_a$, $D(E_af)$ is only guaranteed to be a
compactly supported $H^{-1}$ distribution. Integration by parts in
this distributional sense gives
\[
 \mathfrak g*D(E_af)=i\,\mathfrak g'*E_af=iU_f.
\]
This is the extension on this particular derivative range; it does
not assert a bounded extension of $G_a$ to every $H^{-1}$ distribution.
It retains any boundary contributions of zero extension. Therefore
$\mathcal T_af=0$ is meaningful without making $Df$ an ordinary
$L^2$ vector. Suzuki's shifted derivative-completion and generalized
eigenvalue construction already addresses this distinction
\cite[Sections 8.3--8.5]{Suzuki2026}; the proposition gives its
zero-eigenvalue question as a local continuous integral equation.
It is a derived formulation, not a newly discovered RH equivalence.

\begin{proposition}[A zero eigenvalue does not supply generic $H_0^1$ regularity]
\label{prop:ns51-no-generic-bootstrap}
There is an interval $I_r$, a bounded real self-adjoint rank-one
operator $K$ on $L^2(I_r)$, and a nonzero even
$\tau\in\mathcal D(\tfrac12L_\Delta^{\rm Dir})$ such that
\[
 (\tfrac12L_\Delta^{\rm Dir}+K)\tau=0,
 \qquad \tau\notin H_0^1(I_r).
\]
The operator $K$ is bounded on every $L^p(I_r)$, $1\le p\le\infty$,
and $\tau$ has the continuous zero extension and logarithmic boundary
rate appearing in Theorem~\ref{thm:v118-eigen-endpoint-zero}.
This is a countermodel for inference from the principal operator and
bounded perturbation alone, not a counterexample for the arithmetic
$\mathcal K_\lambda$.
\end{proposition}
\begin{proof}
Choose $0<r<e^{-\gamma_{\rm E}}$. The one-dimensional specialization
of \cite[Theorem 1.2]{HLSLogBoundary} supplies the positive even torsion
function $\tau$ with $L_\Delta^{\rm Dir}\tau=1$ and two-sided bounds
$\tau(x)\asymp\ell(r-|x|)^{1/2}$. Its weak equation with right side
$1\in L^2(I_r)$ puts it in the operator domain. Set
\[
 M=\int_{I_r}\tau(x)\,dx>0,\qquad
 Kf=-\frac{1}{2M}\left(\int_{I_r}f(x)\,dx\right)\mathbf1_{I_r}.
\]
This is bounded, real, self-adjoint and rank one, and $K\tau=-1/2$.
If $\tau\in H_0^1(I_r)$, absolute continuity and Cauchy--Schwarz
would give $|\tau(x)|\le\sqrt{r-|x|}\,\|\tau'\|_2$ near either
endpoint. This contradicts its logarithmic lower bound. The two-sided
boundary estimate also supplies the stated continuous extension.
\end{proof}

\paragraph{Named open input: affine-potential injectivity (API).}
For every finite $a>0$, require
\begin{equation}\label{eq:ns51-api}
 f\in V_a,\quad (\mathfrak g*E_af)|_{I_a}\in
 \operatorname{span}\{1,x\}\quad\Longrightarrow\quad f=0.
\end{equation}
The parity versions must use the constants specified above. No
injectivity estimate for the actual arithmetic kernel is proved here.
Injectivity on all of $L^2(I_a)$ would be sufficient but is a stronger
statement than the one established as equivalent to the physical
kernel condition; no regularity upgrade is being assumed.

\begin{proposition}[Sign continuation with an explicit missing input]
\label{prop:ns51-sign-continuation}
Assume the complete fixed-window realizations above, compact resolvent,
continuity of their lowest eigenvalue in $a$, positivity for sufficiently
small $a$, and API for every finite $a>0$. Then every $A_a$ is strictly
positive at fixed $a$. Consequently the complete compact-test Weil
form is nonnegative and Weil's criterion implies RH. No common positive
spectral gap is required. This is a conditional implication; API is open.
\end{proposition}
\begin{proof}
The first four structural inputs, excluding API, are supplied by the
fixed-window theory and \cite[Theorems 1.3--1.4]{Suzuki2026}.
If the continuous lowest eigenvalue ever became nonpositive, it would
be zero at some finite intermediate window. Compact resolvent makes
that lower edge an eigenvalue. Proposition~\ref{prop:ns51-affine-potential}
then contradicts API. Every smooth compact test lies in a finite window,
so its energy is nonnegative. Apply the classical Weil criterion.
\end{proof}

\paragraph{Scope of the finding.}
The unjustified $H_0^1$ shortcut has been replaced by a valid weak
integral formulation. The generic bootstrap is false, while its truth
for a hypothetical arithmetic nullvector has not been decided. The
remaining obstacle is uniqueness for this specific signed arithmetic
integral equation, not a proved negative direction and not a required
uniform inverse bound. Complete even and odd sectors have both been
retained; no source-complement argument is substituted. Growing
near-zero blocks are consistent with fixed-window injectivity.


% NS-52. Exact finite identities and a conditional obstruction; no RH assumption.
% Standalone fragment: requires amsmath, amssymb, amsthm, hyperref and
% proposition/proof environments. No manuscript version is assigned here.
\subsection*{NS-52: paired heat-zero transport and its first obstruction}

\paragraph{Normalization and claim scope.}
Use the conventional de Bruijn--Newman family
\[
 H_t(z)=\int_0^\infty e^{tu^2}\Phi(u)\cos(zu)\,du,
 \qquad H_0(z)=\tfrac18\xi(\tfrac12+iz/2),
\]
as in Polymath15, equations (1)--(4) and Proposition 3.1
(\url{https://arxiv.org/html/1904.12438#S3}).
Put $Z_t(w)=8H_t(2w)$ and $\kappa=1/4$. Then
\begin{equation}\label{eq:ns52-normalization}
 Z_0(w)=\Xi(w):=\xi(\tfrac12+iw),\qquad
 \partial_t Z_t=-\kappa\partial_w^2 Z_t.
\end{equation}
Thus the heat coefficient in the ordinary zeta-ordinate coordinate is
$1/4$, not $1$. All zero sums below count multiplicity. No assertion
about the locations of the zeros of $Z_0$ is assumed.

For $f,g\in C_c^\infty(-a,a)$ define the nonunitary transform and
its analytic sesquilinear product by
\[
 F_f(w)=\int_{-a}^a f(x)e^{-iwx}\,dx,\qquad
 F_g^*(w)=\overline{F_g(\overline w)},\qquad
 \Psi_{f,g}(w)=F_f(w)F_g^*(w).
\]
The second factor is an entire function of $w$. Replacing it by
$\overline{F_g(w)}$ off the real axis would give the wrong form.
With $\mathcal Z_t$ the zero multiset of $Z_t$, put
\begin{equation}\label{eq:ns52-zero-form}
 q_{a,t}(f,g)=\sum_{w\in\mathcal Z_t}\Psi_{f,g}(w).
\end{equation}
At $t=0$, $\rho=1/2+iw$ and
$\rho^\sharp=1-\overline\rho=1/2+i\overline w$ identify this
exactly with the manuscript's first-slot-linear Weil zero formula:
$q_{a,0}(f,g)=q_a(f,g)$. There is no $2\pi$ factor with the
displayed nonunitary transform. Using the unitary transform instead
would put $2\pi$ in front of the sum. The factor $8$ in $Z_t$ does
not change any zero or its multiplicity. The completed function has
no poles; no additional pole term is to be added to this zero-side
definition. This is not an assertion that the undeformed geometric
prime and pole terms can be heat-evolved separately.

For completeness, the sum in \eqref{eq:ns52-zero-form} is absolutely
convergent on this smooth core for every $t\ge0$, locally uniformly
in $t$. The de Bruijn strip theorem (Polymath15, Theorem 3.2) and
the unconditional initial critical strip give
$|\operatorname{Im} w|\le1/2$. For $0\le t\le T$, Gaussian
majorization of the super-exponentially decreasing $\Phi$ gives
$|Z_t(w)|\le C_T e^{C_T|w|^2}$ and $Z_t(0)\ge Z_0(0)>0$.
Jensen's formula therefore bounds the zero count in $|w|\le R$
by $C_T(1+R)^2$, uniformly in $t$. Repeated integration by parts
gives arbitrary inverse powers of $|\operatorname{Re}w|$ for both
test transforms on the fixed strip. This proves convergence and
uniformly small tails. Local contour integrals count any finite
cluster continuously through collisions, so the complete form is
continuous in $t$ on the smooth core. This argument does not justify
termwise differentiation of the infinite sum or claim a closed
deformed form on the original complete form domain.
If all zeros at a fixed $t_0$ are real, then
\[
 q_{a,t_0}[f]=\sum_{w\in\mathcal Z_{t_0}}|F_f(w)|^2\ge0.
\]
For example the classical de Bruijn theorem supplies such $t_0\ge1/2$.

\begin{proposition}[Finite-cluster transport, including a collision]
\label{prop:ns52-cluster}
Let $Z(t,w)$ be holomorphic near $J\times\overline\Omega$, where
$J$ is an interval of real times and $\Omega$ is a bounded domain
with piecewise smooth positively oriented boundary $\Gamma$.
Assume $Z$ has no zero on $\Gamma$ for $t\in J$ and satisfies
$\partial_t Z=-\kappa Z_{ww}$ with constant $\kappa>0$.
For an entire $\Psi$, define the finite sum $S_\Gamma(t)$ of
$\Psi$ over the zeros in $\Omega$, counted with multiplicity.
Then $S_\Gamma$ is differentiable through all internal collisions and
\begin{align}
 S_\Gamma(t)&=\frac1{2\pi i}\int_\Gamma
                   \Psi(w)\frac{Z_w}{Z}(t,w)\,dw,
                   \label{eq:ns52-contour}\\
 S_\Gamma'(t)&=\frac{\kappa}{2\pi i}\int_\Gamma
                   \Psi'(w)\frac{Z_{ww}}{Z}(t,w)\,dw.
                   \label{eq:ns52-contour-derivative}
\end{align}
At a time with $m$ simple cluster zeros $w_1,\ldots,w_m$, set
$A(w)=Z(t,w)/\prod_{j=1}^m(w-w_j)$ and $b=A'/A$ locally in
the cluster. Then
\begin{equation}\label{eq:ns52-divided-difference}
 S_\Gamma'
 =2\kappa\sum_{1\le j<k\le m}
       \frac{\Psi'(w_j)-\Psi'(w_k)}{w_j-w_k}
   +2\kappa\sum_{j=1}^m b(w_j)\Psi'(w_j).
\end{equation}
If the cluster consists at a particular time of one zero $c$ of
order $m$, write $Z(t,w)=(w-c)^m A(w)$ with $A(c)\ne0$. Its
derivative at that time is exactly
\begin{equation}\label{eq:ns52-multiple}
 S_\Gamma'=
 \kappa m(m-1)\Psi''(c)
       +2\kappa m\frac{A'(c)}{A(c)}\Psi'(c).
\end{equation}
In particular, a double zero contributes
$2\kappa\Psi''(c)+4\kappa(A'/A)(c)\Psi'(c)$.
\end{proposition}
\begin{proof}
The argument principle proves \eqref{eq:ns52-contour}. The nonvanishing
boundary allows differentiation under its compact integral. Since
\[
 \partial_t(Z_w/Z)=-\kappa\partial_w(Z_{ww}/Z),
\]
integration by parts around $\Gamma$ gives
\eqref{eq:ns52-contour-derivative}. At a simple zero,
$w_j'=\kappa Z_{ww}(w_j)/Z_w(w_j)$, and factorization gives
\[
 w_j'=2\kappa\sum_{k\ne j}\frac1{w_j-w_k}
                         +2\kappa b(w_j).
\]
Pair the $j,k$ and $k,j$ contributions to $\sum_jw_j'\Psi'(w_j)$.
This proves \eqref{eq:ns52-divided-difference}. At an order-$m$ zero,
\[
 \frac{Z_{ww}}Z=\frac{m(m-1)}{(w-c)^2}
       +\frac{2m}{w-c}\frac{A'}A+\frac{A''}A.
\]
The residue after multiplying by $\Psi'$ is precisely
\eqref{eq:ns52-multiple}. No labeling or derivative of a colliding
individual root is needed.
\end{proof}

\paragraph{Complete symmetries and the term that must not be omitted.}
For real $t$, $Z_t$ is real entire and even. Choose the union of
cluster contours invariant under conjugation and $w\mapsto-w$,
counting each distinct zero once with its actual multiplicity.
Then the sum defines a Hermitian sesquilinear contribution and the
same identities apply. A double collision at a nonzero real $c$
has a reflected double collision at $-c$; both must be included
in the symmetry-complete contribution. At $c=0$ one would count
one cluster, not duplicate it (the actual $H_t$ has no imaginary-axis
zero). There is no RH assumption here.

For an isolated pair $w_\pm=c\pm d$, the internal term is
$2\kappa[\Psi'(w_+)-\Psi'(w_-)]/(w_+-w_-)$.
It converges to $2\kappa\Psi''(c)$ as $d\to0$.
The two divergent individual velocities therefore cancel exactly.
The remaining $b=A'/A$ term retains the interaction with the rest
of the entire function; it does not vanish by symmetry. The local
analytic factor suffices for the proof, so no additional rearrangement
of an infinite velocity sum is used here. For a full isolated double
cluster the limiting value is the double-zero formula above, not
merely $2\kappa\Psi''(c)$.

\begin{proposition}[The elementary estimate still grows with support]
\label{prop:ns52-support}
For $f,g\in L^2(-a,a)$ and $n\ge0$,
\begin{equation}\label{eq:ns52-pw}
 |\Psi_{f,g}^{(n)}(w)|
 \le 2a(2a)^n e^{2a|\operatorname{Im}w|}\|f\|_2\|g\|_2.
\end{equation}
Consequently, at a double collision with
$|\operatorname{Im}c|\le Y$ and $|(A'/A)(c)|\le B$, the bound
\begin{equation}\label{eq:ns52-support-cost}
 |S_\Gamma'|
 \le16\kappa(a^3+Ba^2)e^{2aY}\|f\|_2\|g\|_2
\end{equation}
holds. The $a^3$ dependence of a single real collision contribution
cannot be replaced by a constant for all compact smooth tests.
Specifically, suppose a double collision occurs at real $c$, and
choose a real odd $h\in C_c^\infty(-1,1)$ with $\|h\|_2=1$ and
$M_1=\int u h(u)\,du\ne0$. For
$f_a(x)=a^{-1/2}h(x/a)e^{icx}$ one has exactly
\begin{equation}\label{eq:ns52-collision-witness}
 \|f_a\|_2=1,\qquad
 S_\Gamma'[f_a]=4\kappa a^3 M_1^2.
\end{equation}
The occurrence of such a collision in the nonnegative-time arithmetic
family is not asserted.
\end{proposition}
\begin{proof}
Cauchy--Schwarz gives
$|F_f^{(j)}(w)|\le a^j\sqrt{2a}e^{a|\operatorname{Im}w|}\|f\|_2$.
Leibniz's formula proves \eqref{eq:ns52-pw}; use $n=1,2$ in
\eqref{eq:ns52-multiple} to obtain \eqref{eq:ns52-support-cost}.
For the witness, $F_{f_a}(c+v)=\sqrt a F_h(av)$, so
$F_{f_a}(c)=0$, $\Psi_{f_a,f_a}'(c)=0$, and
$\Psi_{f_a,f_a}''(c)=2a^3M_1^2$. Apply the double-zero formula.
This is a lower bound on a conditional individual cluster's derivative,
not a lower bound on the complete arithmetic derivative: other clusters
remain present.
\end{proof}

\paragraph{Passage to the whole time-dependent sum.}
Absolute convergence of \eqref{eq:ns52-zero-form} at each time
does not imply convergence of its differentiated cluster series.
One sufficient hypothesis is an exhausting family of finite zero
multisets tracked with multiplicity on $[0,t_0]$, with piecewise
contours as necessary, whose endpoint sums converge to the complete
forms, and whose integrated derivatives satisfy
\[
 \int_0^{t_0}S_n'(t)[f]\,dt\le C\|f\|_2^2
\]
for every compact smooth $f$, with one $C$ independent of the
exhaustion and its support. An integrable bound
$S_n'(t)[f]\le B(t)\|f\|_2^2$, $\int_0^{t_0}B(t)\,dt<\infty$,
would suffice. Full symmetry makes these diagonal quantities real.
These are hypotheses, not consequences of pair cancellation. They
would imply $q_a[f]\ge q_{a,t_0}[f]-C\|f\|_2^2$ by the fundamental
theorem of calculus followed by the endpoint limit. The next result
shows that this particular comparison has an additional obstruction.

\begin{proposition}[An unmatched real zero obstructs uniform
same-test comparison]\label{prop:ns52-unmatched}
Assume the archived cutoff-radical lemma: there exists
$\phi\in\mathcal S(\mathbb R)$ such that, for every fixed finite
linear combination $h$ of its real translates, smooth cutoffs
$\chi_a h\in C_c^\infty(-a,a)$ satisfy
\[
 \chi_a h\longrightarrow h\text{ in }L^2\text{ and }L^1,
 \qquad q_a[\chi_a h]\longrightarrow0.
\]
Assume $q_{a,t_0}$ is the same-test zero form
\eqref{eq:ns52-zero-form}, all its zeros are real, and it has a
zero $c$ of multiplicity $m_c\ge1$ with $F_\phi(c)\ne0$.
Then no finite $C\ge0$ independent of $a$ can satisfy
\begin{equation}\label{eq:ns52-rejected-comparison}
 q_a[f]\ge q_{a,t_0}[f]-C\|f\|_2^2
 \quad(f\in C_c^\infty(-a,a))
\end{equation}
on a cofinal family of windows.
\end{proposition}
\begin{proof}
Fix $T>0$. The Schwartz assumption implies
\[
 B_T:=\sum_{k\in\mathbb Z}
       |\langle\phi,\phi(\cdot-kT)\rangle|<\infty.
\]
For each fixed positive integer $M$ put
\[
 h_M=M^{-1/2}\sum_{j=1}^M e^{icjT}\phi(\cdot-jT).
\]
Counting equal differences of shifts gives
$\|h_M\|_2^2\le B_T$, whereas
$F_{h_M}(c)=\sqrt M F_\phi(c)$.
Since all deformed zeros are real,
$q_{a,t_0}[f]\ge m_c|F_f(c)|^2$ for every smooth test.
Apply \eqref{eq:ns52-rejected-comparison} to $\chi_a h_M$ and
let $a\to\infty$ along the prescribed cofinal family, keeping $M$
fixed. The $L^1$ convergence controls the single real Fourier
evaluation. Thus
\[
 m_c M|F_\phi(c)|^2\le C\|h_M\|_2^2\le CB_T.
\]
Now let $M\to\infty$, a contradiction. No cutoff estimate uniform
in $M$, global closedness, or lower semicontinuity of a deformed
form is used.
\end{proof}

The cutoff input is Lemma~\ref{lem:v136-total-radicals}: its ordinary
norm convergence and vanishing complete operator residual imply the
displayed form limit. Schwartz decay and dominated convergence supply
the additional $L^1$ limit. The source used there has Fourier transform
a nonzero constant times $\Xi$, by \eqref{eq:v150-mellin} and the
source normalization, so $F_\phi(c)\ne0$ means precisely that this
real deformed zero is not an original Xi zero. An unmatched zero
at a specified $t_0$ has not been certified or proved in this
checkpoint. The proposition states its hypothesis explicitly; it
does not assert an unconditional arithmetic closure. It also does
not assert any negative direction of the original Weil form.

\paragraph{Disposition.}
The finite algebraic test succeeds: multiplicities and complete
pairing remove the inverse-gap singularity. What had to change is
the normalization, the analytic off-real sesquilinear product, and
retention of the exterior-cluster logarithmic derivative. What is
still missing is a complete signed transport estimate with support
independent error. Moreover the proposed same-test positive-time
comparison is impossible whenever there is an unmatched real
deformed zero. A broader heat approach would therefore need a
different comparison or a justified transport of tests that respects
the original radical; no such transport is supplied. These are
finite exact identities, proved implications with explicit hypotheses,
and a named open transport input. No new window, tail metric,
numerical certificate, G2 assertion, or RH assertion is made.


% NS-53. Exact identities, proved implications, and an explicitly open input.
% No finite numerical approximation is offered as convergence evidence.
\subsection{NS-53: residualized integer-dilation blocks}
\label{sec:ns53-nb-blocks}

Work in the real Hilbert space $H=L^2((0,\infty),dx)$, with
$\chi=\mathbf 1_{(0,1)}$ and $\rho_n(x)=\{1/(nx)\}$ for integers $n\geq1$.
Real coefficients lose nothing for this real target; the complex version
uses conjugate transposes. Put
\[
 V_N=\operatorname{span}\{\rho_1,\ldots,\rho_N\},\qquad
 r_N=\chi-P_N\chi,\qquad E_N=\|r_N\|^2=d_N^2,
\]
where $P_N$ is the orthogonal projection onto $V_N$.
The strong Nyman--Beurling criterion is the established equivalence
$E_N\to0\Longleftrightarrow\mathrm{RH}$; see B\'aez-Duarte,
\href{https://arxiv.org/abs/math/0202141}{Theorem~1.1}.
No particular upper rate is required. Burnol's
\href{https://doi.org/10.1006/aima.2001.2066}{lower-bound theorem}
does not require a matching upper bound. This checkpoint proves no
convergence statement and makes no G2 or RH claim.

\begin{proposition}[Exact arithmetic entries and finite independence]
\label{prop:ns53-nb-entries}
Let $G_{mn}=\langle\rho_m,\rho_n\rangle$, $b_n=\langle\chi,\rho_n\rangle$,
and $\kappa=\log(2\pi)-\gamma$, where $\gamma$ is Euler's constant.
Then
\begin{align}
 G_{mn}&=\int_0^\infty\{y/m\}\{y/n\}\,\frac{dy}{y^2},
 & b_n&=\frac{\log n+1-\gamma}{n},
 &G_{nn}&=\frac{\kappa}{n}.\label{eq:ns53-nb-entries}
\end{align}
In particular, writing $a_j=\lfloor j/m\rfloor$ and
$d_j=\lfloor j/n\rfloor$, the following is an exact convergent series:
\begin{equation}
 G_{mn}=\frac1{mn}+\sum_{j=1}^{\infty}
 \left[\frac1{mn}-\left(\frac{d_j}{m}+\frac{a_j}{n}\right)
 \log\left(1+\frac1j\right)+\frac{a_jd_j}{j(j+1)}\right].
 \label{eq:ns53-nb-gram-series}
\end{equation}
Each bracket is a nonnegative cell integral. Every Gram matrix belonging
to finitely many distinct positive integers is positive definite.
Moreover $E_N>0$ for every finite $N$.
\end{proposition}

\begin{proof}
Substitute $y=1/x$ and integrate on $[j,j+1)$; both floor functions
are constant on that interval. This gives the integral and series formulas.
The change $y=nt$ gives the diagonal scaling and
\[
 b_n=\frac1n\left(\log n+\int_1^\infty\frac{\{t\}}{t^2}\,dt\right).
\]
The integral truncated at an integer $K$ is $1+\log K-H_K$.
Similarly, integrating $\{t\}^2/t^2$ on the first $M+1$ unit cells gives
\[
 2M+2-H_{M+1}-2M\log(M+1)+2\log(M!).
\]
Stirling's formula gives its limit $\kappa$. For finite independence,
a vanishing linear combination of $\{y/n\}$ is piecewise affine.
At the smallest index with nonzero coefficient its jump cannot be
cancelled by a smaller nonzero coefficient, a contradiction. If
$\chi=\sum_{n\leq N}c_n\rho_n$, comparison of the jumps at every positive
integer in the $y$ variable gives
$\sum_{n\mid j}c_n=-\mathbf 1_{j=1}$, with $c_n=0$ for $n>N$.
M\"obius inversion would force $c_n=-\mu(n)$ for every $n$, contradicted
by a prime greater than $N$. A finite-dimensional span is closed, so
$E_N>0$.
\end{proof}

\begin{proposition}[Exact block improvement, including all cross terms]
\label{prop:ns53-nb-block-gain}
Fix integers $1\leq N<M$. Let $G$ be the $N$ by $N$ old Gram matrix,
$C=(G_{mn})_{1\leq m\leq N,\,N<n\leq M}$, and
$D=(G_{mn})_{N<m,n\leq M}$.
Write $b=(b_m)_{m\leq N}$, $b_B=(b_n)_{N<n\leq M}$, and define
\begin{equation}
 c=G^{-1}b,\qquad h=b_B-C^{\mathsf T}c,\qquad
 S=D-C^{\mathsf T}G^{-1}C.
 \label{eq:ns53-nb-schur}
\end{equation}
Then $S$ is positive definite, and
\begin{equation}
 E_N-E_M=h^{\mathsf T}S^{-1}h
 =\sup_{w\ne0}\frac{|w^{\mathsf T}h|^2}{w^{\mathsf T}Sw}.
 \label{eq:ns53-nb-gain}
\end{equation}
The gain is zero if and only if $h=0$; strict finite improvement is not
asserted for every block. The coefficients for the enlarged optimum are
$(c-G^{-1}CS^{-1}h,S^{-1}h)$. In addition,
\begin{equation}
 E_N-E_M\ \geq\
 \frac{\sum_{n=N+1}^M|b_n-\sum_{m\leq N}c_mG_{mn}|^2}
 {\kappa\sum_{n=N+1}^M1/n}
 \ \geq\
 \frac{\|h\|_2^2}{\kappa\log(M/N)}.
 \label{eq:ns53-nb-trace-gain}
\end{equation}
These are finite, unconditional identities and inequalities, not estimates
of the numerator as $N$ grows.
\end{proposition}

\begin{proof}
Set $z_n=(I-P_N)\rho_n$ for $N<n\leq M$. Its Gram matrix is $S$,
and $\langle r_N,z_n\rangle=h_n$. Finite independence shows that these
$z_n$ are independent. The orthogonal decomposition
$V_M=V_N\mathbin{\oplus}\operatorname{span}\{z_n\}$ gives
\eqref{eq:ns53-nb-gain} and the coefficient update. Since
$0<S\leq D$, the largest eigenvalue of $S$ is at most
$\operatorname{tr}D=\kappa\sum_{n=N+1}^M1/n$. Therefore
$h^{\mathsf T}S^{-1}h\geq\|h\|_2^2/\operatorname{tr}D$.
The decreasing function $1/x$ gives
$\sum_{n=N+1}^M1/n\leq\log(M/N)$.
\end{proof}

\begin{remark}[Conditioning is retained]
Replacing $S$ by $D$ in the exact identity, or using $b_B$ instead of $h$,
would discard the old--new interaction. A small eigenvalue of $S$ can
make the exact gain much larger than the trace lower bound; it causes no
missing inverse error in the analytic lower bound. For a proposed real
block vector $w$, its exact gain after optimizing its scalar coefficient
is $|w^{\mathsf T}h|^2/(w^{\mathsf T}Sw)$. If
$\sum_{n=N+1}^M w_n/n=0$, the unprojected combination
$\sum w_n\rho_n$ vanishes on $x>1/(N+1)$. Its residualized combination
need not have this support, and its norm is still $w^{\mathsf T}Sw$.
\end{remark}

\begin{proposition}[Exact divisor cells and the raw M\"obius-prefix penalty]
\label{prop:ns53-nb-divisor-cells}
For arbitrary real coefficients $c_1,\ldots,c_N$, put
\[
 A=\sum_{n\leq N}\frac{c_n}{n},\qquad
 B_j=1+\sum_{n\leq N}c_n\left\lfloor\frac jn\right\rfloor
     =1+\sum_{k=1}^j\ \sum_{\substack{n\mid k\\n\leq N}}c_n.
\]
For $r=\chi-\sum c_n\rho_n$, one has $r(1/y)=-Ay$ on $(0,1)$
and $r(1/y)=B_j-Ay$ on $[j,j+1)$, $j\geq1$. Consequently
\begin{equation}
 \|r\|^2=A^2+\sum_{j=1}^{\infty}
 \left[A^2-2AB_j\log\left(1+\frac1j\right)
          +\frac{B_j^2}{j(j+1)}\right].
 \label{eq:ns53-nb-divisor-energy}
\end{equation}
The brackets are nonnegative cell integrals; they are not to be split
into three separately summed series. For the raw choice $c_n=-\mu(n)$,
\begin{equation}
 \left\|\chi+\sum_{n\leq N}\mu(n)\rho_n\right\|^2
 \geq (N+1)\left|\sum_{n\leq N}\frac{\mu(n)}n\right|^2.
 \label{eq:ns53-nb-mobius-penalty}
\end{equation}
\end{proposition}

\begin{proof}
Expand each fractional part. The energy identity is the integral of the
squared affine function on each unit cell. For the M\"obius choice,
$B_j=0$ for $1\leq j\leq N$, by
$\sum_{n\mid k}\mu(n)=\mathbf 1_{k=1}$. The interval $(0,N+1)$ alone
then contributes the displayed lower bound.
\end{proof}

\begin{remark}[What the arithmetic lemma does and does not supply]
Equation~\eqref{eq:ns53-nb-mobius-penalty} gives an explicit necessary
condition for this particular unsmoothed coefficient choice to converge:
$\sqrt N\sum_{n\leq N}\mu(n)/n\to0$. It is not a lower bound for the
optimal $E_N$. The divergence of the raw natural approximation is already
recorded in B\'aez-Duarte's Introduction; no novelty is claimed for its
failure. Smoothing or adaptive coefficients change the problem and must
be evaluated through the full residual and Gram matrix.
\end{remark}

\begin{proposition}[A sufficient block-correlation input at any rate]
\label{prop:ns53-nb-any-rate}
Fix an integer $q\geq2$, $N_j=q^j$, and real numbers
$0\leq\alpha_j<1$ for $j\geq j_0$, with $\sum_j\alpha_j=\infty$.
Use the actual optimal coefficients $c^{(N_j)}=G_{N_j}^{-1}b_{N_j}$.
Suppose the following \emph{residual block-correlation input} holds:
\begin{equation}
 \sum_{n=N_j+1}^{qN_j}
 \left|\frac{\log n+1-\gamma}{n}
       -\sum_{m\leq N_j}c_m^{(N_j)}G_{mn}\right|^2
 \geq\kappa\log q\,\alpha_j E_{N_j}.
 \tag{RBC}\label{eq:ns53-nb-rbc}
\end{equation}
Then $E_N\to0$. For example, the hypothetical bound (RBC) with
$\alpha_j=1/(2(j+1))$ would give
$E_{q^j}=O(j^{-1/2})$, hence $E_N=O((\log N)^{-1/2})$.
This slower-than-optimal rate would already meet the established criterion.
\end{proposition}

\begin{proof}
Equation~\eqref{eq:ns53-nb-trace-gain} gives
$E_{N_{j+1}}\leq(1-\alpha_j)E_{N_j}$. Iteration and
$1-t\leq e^{-t}$ prove convergence. Monotonicity supplies intermediate
$N$. For the example, summing $1/(2(j+1))$ yields the stated rate.
\end{proof}

\paragraph{Named open input and disposition.}
(RBC) is \emph{unproved}; it is a specific sufficient arithmetic estimate,
not asserted equivalent to RH or known to follow from RH. It may be
stronger than needed because the trace bound can lose the small
Gram eigenvalues that amplify the exact gain. The weaker exact target is
a nonsummable sequence of relative gains
$h_j^{\mathsf T}S_j^{-1}h_j/E_{N_j}$, which is equivalent to convergence
and is not itself a new theorem. The arithmetic work still missing is
a lower estimate for the corrected correlations, or for an explicitly
constructed block quotient, for the \emph{actual optimal residuals}.
Entry formulas, a finite collection of positive gains, and a finite list of distances
do not supply it. No such asymptotic gain estimate is proved here.
The missing step is the asymptotic gain estimate; this establishes
neither failure of (RBC) nor failure of the approximation object. No new Weil window, tail metric, G2 or RH
claim is introduced.
